\documentclass[11pt]{article}
\usepackage[T1]{fontenc}
\usepackage{lmodern,microtype}
\usepackage{amsmath,amssymb,booktabs,enumitem}
\usepackage[letterpaper,margin=1in]{geometry}
\usepackage[colorlinks=true,urlcolor=blue,citecolor=blue]{hyperref}
\newcommand{\Q}{\mathbb Q}
\newcommand{\F}{\mathbb F}
\newcommand{\Z}{\mathbb Z}
\title{Metaplectic Jacquet--Langlands Transfer, Companion Paper (Rank Two, Wild Place)\\Revision Log: Draft 1 to Draft 2}
\author{Prepared by Claude (Anthropic) for David Joyner}
\date{November 8, 2026}
\begin{document}
\maketitle

\section*{Scope}
Draft 2 responds to ChatGPT's referee report on Draft 1 (recommendation:
major revision).  The report accepted the finite Fourier calculations
after independent reconstruction, but did not accept Target A as proved,
for two reasons concerning Proposition \texttt{prop:H1}: the written
proof of case (i) contained an incorrect assertion ($\beta(X,X')=0$
whenever $v(t)>0$), exhibited by the explicit pair
$X=1+\varpi^3j$, $X'=1+\varpi^3+\varpi^4j$ with $\beta=2$; and cases
(iv)--(v) rested on a stratified sample rather than a proof.  The
referee asked that (H1) be repaired without new hypotheses, that the
proof of Target A then be reexamined for remaining dependencies, that the
finite Fourier work be preserved, and that the Python modules be
supplied.

\section*{The repair of (H1)}
The referee's example is correct: in the explicit model the pair above
has $\beta=2$, $\kappa(X)=0$, $\kappa(X')=1$, $\kappa(XX')=0$, so the
identity holds but the Draft 1 argument for it was wrong.  The error was
the claim that $1+\varepsilon$ in $r=(1+t)(1+\varepsilon)$ lies in
$U^{(4)}$; it lies only in $U^{(3)}$, and $(U^{(3)},\varpi)_{E_0}\neq1$.

Draft 2 replaces the case split entirely.  Two changes make a complete
hand proof possible.
\begin{enumerate}[leftmargin=*]
\item \textbf{A cleaner closed form for the star section.}  For
$X=x+yj\in K^6$ write $X=zv^3$ with $z\in U_F^{(1)}$, $v\in K^2$.  Taking
reduced norms gives $z\equiv\mathrm{Nrd}(X)^{-1}$ modulo $F^{\times3}$, so
the star lift is $\ell(X)=(X,\kappa(X))$ with
\[
 \kappa(X)=(y,\mathrm{Nrd}\,X)_{E_0}\quad(y\neq0),\qquad \kappa(X)=0\quad(y=0),
\]
(Corollary \texttt{cor:starK6}).  This is equivalent to the Draft 1 digit
formula $c(X)\,v(y)$ but needs no digits, and it is what makes the
identity below transparent.
\item \textbf{The right factorization.}  In the main case
($y,y'\neq0$, $XX'\notin E_0$), with $\tau=\sigma(y'/y)$ and
$R=x'+\sigma(x)\tau$, the identity (H1) is equivalent to
\[
 (R,\tau)_{E_0}=v(\tau)\lambda(x')-v(R)\bigl(\lambda(x')-\lambda(x)\bigr),\qquad
 \lambda(x):=(x,\varpi)_{E_0}.
\]
Factoring $R=x'(1+\tau_2)$ with $\tau_2=\sigma(x)\tau/x'$ and writing
$\tau=\tau_2\cdot x'/\sigma(x)$, the left side splits into three
symbols: $(x',\tau)$, which is $v(\tau)\lambda(x')$ because $x'\in U^{(3)}$
kills $\mu_8$ and $U^{(1)}$; $(1+\tau_2,\tau_2)=0$ by Steinberg; and
$(1+\tau_2,x'/\sigma x)$, which is $v(R)$ times
$(\varpi,x')-(\varpi,\sigma x)$ by the same depth bound and Galois
invariance.  No case distinction on $v(\tau)$ or on $v(1+\tau)$ occurs;
the former boundary cases (iv) are not special.  The degenerate cases
($y=0$, $y'=0$, both, or $XX'\in E_0$) are each a few lines of the same
kind.  The only inputs are Kubota's cocycle formula, bimultiplicativity,
antisymmetry, Steinberg, Galois invariance, triviality on cubes, and the
depth bound $(U^{(i)},U^{(j)})_{E_0}=1$ for $i+j\geq4$ (Fesenko--Vostokov;
Serre).  No new hypothesis is introduced, and the explicit values of the
symbols are not used in this proof.
\end{enumerate}
Proposition \texttt{prop:H1} is now labeled [P].  A new script
(\texttt{step12\_H1\_proof\_check.py}) reproduces the referee's example,
checks the new closed form against the digit form on $3000$ elements,
checks the symbol facts used, and checks the reduced identity and (H1)
on $12{,}000$ structured instances covering every case of the proof
($\beta\neq0$ in about $44\%$); this is corroboration only.

\section*{Reexamination of Target A: dependency ledger}
A new Remark \texttt{rem:dependencies} in Section 11 itemizes everything
the main theorem uses:
\begin{itemize}[leftmargin=*]
\item[(D1)] [E] Part I (Definition 2.2, Propositions 2.3, 3.1, 3.5(a),
Lemma 3.4, Theorem 3.9, Remark 8.1) and Part III (normalization of
$L_D$, $b_D=3$, $c_D=27$, $R=1$ on $2$-power elements).
\item[(D2)] [E] Kubota's cocycle on $\mathrm{GL}_2(E_0)$ and the standard
properties of the cubic Hilbert symbols.
\item[(D3)] [P$^c$] the explicit symbol tables (exact linear algebra via
Moore's theorem), used only through $(1+\varpi^3,\varpi)\neq1$ in the
exclusion of $e=0$.
\item[(D4)] [P$^c$] four exhaustive finite enumerations: $\bar G_3$ from
the $648$ depth-one classes; the cube map on $Q/Q_7$; the actions on
$V$; the design search with the Heisenberg extension and Fourier
identities.  The referee reconstructed the design data independently.
\item[(D5)] [P] everything else, including (H1).
\end{itemize}
No sampled computation is load bearing.  The Lazard-splitting route of
Draft 1 is retained as a remark (second route), no longer needed.

\section*{Status changes}
\begin{center}
\begin{tabular}{lll}
\toprule
Item & Draft 1 & Draft 2\\
\midrule
Lemma \texttt{lem:model} (the model is the cover) & [P$^c$] & [E]/[P] (cocycle cited, not sampled)\\
Corollary \texttt{cor:starK6} & [P] & [P] (new closed form)\\
Proposition \texttt{prop:H1} & [P$^c$], case (i) wrong & [P]\\
Target A & [P$^c$] with one sampled step & [P$^c$], all computations exhaustive\\
Level-zero rows and norm twists & [P$^c$] & [P$^c$]\\
Target B positive level; Target C & [O]/[C] & [O]/[C]\\
\bottomrule
\end{tabular}
\end{center}
The label [P$^c$] on Target A now means exactly (D3)--(D4): finite
enumerations over explicitly described sets, each with a supplied exact
script.  The definition of the label in Section 1 was tightened
accordingly (``exhaustive'' is now part of it).

\section*{Items from the referee report not yet addressed}
The report's table marks Sections 7 and 8 ``promising; some arguments
require fuller justification'' and ``minor gaps to repair'' without
specifics in the portion received.  The measure-theoretic step in
Lemma \texttt{lem:cube}(f) (cubing is a local homeomorphism near regular
elements) is the place in Section 7 most likely meant; it is
corroborated by script 2 in $Q/Q_7$ but could be replaced by a direct
Hensel argument.  Section 8 was not changed in Draft 2.  Both are
candidates for Draft 3 once the referee's specifics are available.

\section*{Editorial changes}
Abstract and Section 1 updated; status table splits the Section 10 row
into the model lemma ([E]/[P]) and the proofs ([P]); Remark
\texttt{rem:H1-status} removed and replaced by Remark
\texttt{rem:H1-checks}; Remark \texttt{rem:dependencies} added;
Appendix B lists script 12; bibliography adds Kubota, Fesenko--Vostokov,
Serre, Symonds--Weigel.  The acknowledgments record that the Draft 1
error was found by ChatGPT's report.  Sections 2--9 and the finite
Fourier material are unchanged, as the referee requested.  American
English spelling throughout.

\section*{Deliverables}
\texttt{mjl\_rank2\_wild\_draft2.tex/.pdf} (27 pages), this log, and
\texttt{mjl\_rank2\_wild\_scripts.tar.gz} containing \texttt{quat.py},
\texttt{cover.py} (exact $\Q(\zeta_{12})$ arithmetic; the cubic Hilbert
symbols are solved in \texttt{step7\_symbols.py} from the Steinberg
relations and stored in \texttt{symbols.pkl}), and
\texttt{step1}--\texttt{step12}.  To verify the cocycle and symbol
implementation independently, run \texttt{step7\_symbols.py} first
(regenerates \texttt{symbols.pkl}), then \texttt{step12\_H1\_proof\_check.py}.
\end{document}
